\documentclass[a4paper,12pt]{article}
\usepackage{bib/trymtex}

\addbibresource{bib/references.bib}
\graphicspath{{fig/}}

\begin{document}
\begin{titlepage}
	\newcommand{\HRule}{\rule{\linewidth}{0.5mm}}
	\begin{tikzpicture}[remember picture, overlay]
		% NTNU logo
		\node[anchor=north west, xshift=1.0cm, yshift=-1.0cm] at (current page.north west) {
			\includegraphics[width=2.0cm]{bib/ntnu_logo_liten.png}
		};
	\end{tikzpicture}

	\center

	\vspace*{1cm}

	% Course code & title
	{\color{ntnu-blue}\sffamily\Huge\bfseries TIØ4146 \par}
	\vspace{0.5cm}
	{\sffamily\huge\bfseries Teknologiledelse \par}

	\vspace{1.5cm}

	\HRule
	\vspace{0.5cm}

	% Assignment title
	{\Large\sffamily Gruppeoppgave 2 - Organisasjonsteori\par}
	\vspace{0.3cm}
	%{\Huge\sffamily\textit{EcoBaseNO}\par}

	\vspace{0.5cm}
	\HRule

	\vspace{2cm}

	% Author info
	\begin{minipage}{0.6\textwidth}
		\begin{flushleft}
			\large
			\textbf{Gruppe 186}\\
			10059 \\
			10645 \\
			10813 \\
			10631 \\
			10323 \\
			
		\end{flushleft}
	\end{minipage}%
	\begin{minipage}{0.4\textwidth}
		\begin{flushright}
			\large
			\textbf{Semester:}\\
			Vår 2025
		\end{flushright}
	\end{minipage}

	\vfill

	% University logo/name
	\begin{center}
		{\color{ntnu-blue}\sffamily\Large Norges teknisk-naturvitenskapelige universitet}\\
		\vspace{0.3cm}
		{\sffamily\large Institutt for industriell økonomi og teknologiledelse}

		\vspace{0.5cm}
		{\large 27. Mars 2025}
	\end{center}

	\vspace{1cm}
\end{titlepage}




\setcounter{page}{1}
\renewcommand{\contentsname}{Innholdsfortegnelse}
\tableofcontents
\newpage

\section*{Innledning}

Dette prosjektet er delt inn i tre hovedoppgaver, der alle tre er knyttet til kapitlene 6-10 i pensumboken Teknologiledelse for ingeniørstudier, 2.utgave. 
Oppgavene tar for seg ulike sider ved teknologiledelse og fokuserer på både teori og anvendelse i praksis.

I første hovedoppgave ser på vi på Ski-VM i Trondheim som et prosjekt, der vi skal analysere prosjektet med fokus på blant annet organisering og prosjektledelse. 

Videre, i oppgave 2, er fokuset arbeidslivet og personlig motivasjon for valg av arbeidsgiver. Her skal vi se på en fremtidig arbeidsplass og bruke teori fra pensumboken sammen med offentlig informasjon for å forklare hva som gjør denne bedriften attraktiv. 

Den tredje og siste oppgaven tar for seg to aspekter ved den norske arbeidslivsmodellen som bør ivaretas, samt utfordringer i norsk arbeidsliv knyttet til Gen Z, dagens generasjon. Oppgaven bygger på aktuell debatt og teori fra kapittel 10 i pensum.

\section{Oppgave 1 - Ski-VM i Trondheim}
\subsection{Oppbygning og organisering av Ski-VM}
Amund

\subsection{Var Ski-VM en suksess?}
Det finnes ingen konkret, universalt svar på om Ski-VM i Trondheim er en suksess eller ikke.  
I likhet med det meste annet er dette subjektivt, og dermed er det viktigere spørsmålet for hvem 
var det en suksess og for hvem var det ikke? For et prosjekt med størrelsen av Ski-VM er det 
enormt med aktører og interessenter, derfor velger vi tre som er spesielt interessante, 
Trondheim kommune, Næringslivet rettet mot Ski-VM og Skiforbundet. Dette er aktører som hadde 
spesielt mye innflytelse på-, og var spesielt påvirket av resultatet og konsekvensene av 
arrangementet. For å analysere disse gruppene benyttes en PESTEL-analyse\cite{Teknologiledelse}(Mangler side).

\subsubsection{Trondheim kommune}
Trondheim kommune var den viktigste interessenten av Ski-VM da de klart hadde mest 
innflytelse på Ski-VM. Dette er en stor og kompleks gruppe med mange ulike seksjoner som 
alle har ulike mål og ønsker til Ski-VM, likevel kan kommunens ønsker inndeles i 
følgende hovedmål: Økonomisk bærekraft, politisk gevinst og Miljømessig bærekraft og troverdighet.

\paragraph{Økonomisk bærekraft:} Trondheim kommune investerte tungt i Granåsen med omlag 1.2 
milliarder kr satt inn i prosjektet\cite{NRK12Milliarder}. Disse enorme summene har blitt møtt med skepsis,
og i ettertid er det ennå ikke tydelig om VM har direkte gått i pluss for Trondheim kommune\cite{AftenpostenNyVM}. 
Indirekte kan arrangementet ha vært økonomisk bærekraftig gjennom å stimulere den lokal 
økonomien, det er tydelig at det har vært en enormt økt aktivitet i hotell- og 
restaurantnæringen\cite{DagbladetPriser}. Kommunen har siden signalisert at fremtidige mesterskap ikke
kan få like mye finansiering fra Trondheim\cite{NeaRadioVM}, noe som kan tyde på at VM har hatt mindre økonomisk 
gevinst enn forventet. Samtidig forblir mye av investeringen igjen til gevist for Trondheim i form 
et oppgradert, varig skianlegg.

\paragraph{Politisk gevinst:} VM styrket Trondheims nasjonale omdømme. Mer enn 519 milioner så på 
Ski-VM\cite{AftenpostenSeertall} og noen dager var det mer enn 40 tusen publikummere i 
Granåsen\cite{NRKFolkefest}. Statsministerens støtte fremhevet kommunens 
gjennomføringsevne\cite{Trondheim2025Midler}. 
Arrangementet ble omtalt som mer et skimesterskap, med særlig vekt på frivillighet og 
inkludering\cite{Trondheim2025Baerekraft}.  Dette samsvarer med kommunens politikk for sosial inkludering, 
og bidro til økt lokal stolthet og tillit til kommunens prosjektkapasitet. Undersøkelser viser
at majoriteten av de som deltok i undersøkelsen "hadde et positivt totalinntrykk av VM og 69 prosent 
av dem i aldersgruppen 15-29 år syntes regionen er mer attraktiv etter mesterskapet"\cite{AftenpostenNyVM}.

\paragraph{Miljømessig troverdighet:} Kommunen og arrangørene hadde mål om et klimanøytralt 
mesterskap\cite{TrondheimKommuneVM}. Blant tiltakene var bruk av utslippsfri energi, grønn 
transport og massiv avfallssortering. Klimagassutslipp ble redusert med 98\% sammenlignet 
med Oslo-VM 2011\cite{Trondheim2025Kutt}. Granåsen ble bygget med miljøhensyn og skal ha lang levetid. 

Kommunens krav presset frem innovative løsninger og styrket Trondheims
miljøprofil nasjonalt og internasjonalt.

\subsubsection{Næringslivet rettet mot Ski-VM}
Næringslivet rettet mot Ski-VM var også en viktig interessent. Denne gruppen består av 
hoteller, restauranter, transporttjenester og andre virksomheter som i stor grad hadde sitt 
daglige virke koblet til Ski-VM, enten direkte gjennom tjenester til publikum og deltakere, 
eller indirekte gjennom økt aktivitet i byen. Spesielt viktig innad i denne gruppen er 
næringslivet som ble skapt utelukkende for å ta fordel av Ski-VM, da matbodene og VM-lavvoen 
i sentrum. For næringslivet var hovedmålet å oppnå økt omsetning, bygge nettverk, og 
styrke Trondheim som en attraktiv destinasjon for fremtidige arrangementer og turisme. (Kilde?)

\paragraph{Turisme og markedsføring:} VM utløste økt turisttrafikk og høy hotellbelegg, nær 80\,\%,
til tross for dobling i rompriser\cite{Finansavisen2025}. Internasjonal TV-dekning 
markedsførte Trondheim globalt. Gjennom SPOR-nettverket og samarbeid på tvers av kommuner 
ble regionen profilert som reisemål\cite{NiT2024}. Bedrifter fikk synlighet og brukte VM 
som markedsføringsarena, både nasjonalt og internasjonalt.

\paragraph{Lønnsomhet og lokal verdiskaping:} Arrangementet tiltrakk over 300\,000 
besøkende\cite{Norengros2024}, og næringslivet meldte om vesentlig økt omsetning. 
VM-selskapet benyttet lokale leverandører til alt fra mat og utstyr til miljøtjenester. 
Eksempler inkluderer Norengros og samarbeidsprosjektet “Gauldalstunet”, hvor seks 
nabokommuner presenterte lokalmat\cite{NiT2024}. Mesterskapet skapte nye 
forretningsforbindelser og ga varige ringvirkninger.

\paragraph{Arbeidskraft og kompetanse:} Rundt 2400 frivillige stilte 
opp\cite{TrondheimKommune2023}, og mange bedrifter ansatte ekstra personale for å møte 
etterspørselen. Dette gav arbeidserfaring, spesielt til unge. Næringslivet og kommunen 
samarbeidet om inkludering av personer utenfor arbeidslivet, noe som styrket sosial bærekraft. 
Erfaringene fra logistikken rundt arrangementet styrket den profesjonelle kapasiteten 
i flere bransjer, og teknologiselskaper brukte VM som testarena for nye løsninger. 
Dette kan åpne nye markedsmuligheter og styrke fremtidig næringsutvikling.


\subsubsection{Norges Skiforbund som interessent}
Skiforbundet er en annen sentral interessent og har hatt stor påvirkning på Ski-VMs utforming 
og gjennomføring. Som hovedeier (60 \%) i arrangørselskapet Ski-VM Trondheim 2025 AS og øverste 
sportslige myndighet hadde Norges Skiforbund (NSF) en todelt rolle: strategisk prosjekteier med 
ansvar for nasjonale mål -- rekruttering, prestasjoner og omdømme -- og kommersiell aktør som 
skulle sikre et økonomisk overskudd til videre drift og talentutvikling\cite{ProffSkiVM2025}.

\paragraph{Hovedinteresser}

\begin{itemize}
    \item \textbf{Økonomi} -- overskudd som kan snu flere år med røde tall og 
    medlemstap\cite{Adresseavisen}

    \item \textbf{Sportslig prestasjon} -- norske medaljer som styrker merkevaren Langrenn Norge.
    
    \item \textbf{Bærekraftig arrangement} som viser forbundets samfunnsansvar \cite{TrondheimKommune}.

    \item \textbf{Omdømmebygging} mot frivillige, ungdom og sponsorer -- VM som en del av f
    orbundets rekrutteringsstrategi \cite{OsloVM}.
\end{itemize}

\paragraph{PESTEL-analyse}

\begin{table}[ht]
    \centering
    \begin{tabular}{@{}p{2.7cm}p{10.2cm}@{}}
        \toprule
        \textbf{Faktor}        & \textbf{Relevans for NSF}                                                                                                                                                  \\ \midrule
        \textbf{P}olitiske     & Statlig og kommunal finansiering avhang av at arrangementet leverte brede samfunnsgevinster \cite{TrondheimKommune}.                                                       \\
        \textbf{E}konomiske    & 93~\% av billettene ble solgt nasjonalt; samtidig presset valutarisiko og økte anleggskostnader resultatmarginen \cite{AdressaKjopefest}.                                  \\
        \textbf{S}osiale       & 230\,000 stadiontilskuere og 3\,800 frivillige skapte folkefest, men synkende klubbmedlemskap viser at grasroteffekten ikke er realisert ennå \cite{Adresseavisen,OsloVM}. \\
        \textbf{T}eknologiske  & Prosjektet \textit{Snow for the Future} med automatisert og energieffektiv snøproduksjon gav et eksportbart kompetansefortrinn\cite{Trondheim2025Sustainability}.          \\
        \textbf{E}nvironmental & Klimatiltak som kortreist mat og 10\,000 nyplantede trær reduserte kritikk, men flyreiser fra publikum forble et omdømmerisiko.                                            \\
        \textbf{L}ovmessige    & FIS-reglement, antidopingkrav og offentlige anskaffelsesregler krevde profesjonell prosjektstyring og åpenhet.                                                             \\ \bottomrule
    \end{tabular}
\end{table}

\paragraph{Påvirkning og resultater}
Fra et interessent-perspektiv (pensum kap. 9) er suksess mer enn «on time/on budget»; en eier må 
oppleve \textit{impact success} -- varige gevinster. Ski-VM leverte rekordpublikum, høy TV-rating 
og norsk medaljedominans, noe som styrket NSF både sportslig og omdømmemessig.
Bærekraftsprofilen traff unge målgrupper og kan hjelpe forbundet i kampen om frivillige 
og sponsorer.
Økonomisk ble det rapportert et moderat pluss -- lavere enn håpet, men viktige inntekter vil 
fortsatt komme fra etter-arrangement og markedsrettigheter \cite{Trondheim2025Portal,Adresseavisen}.

\paragraph{Suksessvurdering}
Med tydelige gevinster på tre av fire hovedinteresser kan Ski-VM regnes som en betinget suksess for NSF. Den største svakheten er at medlemstallene foreløpig ikke viser vekst -- et signal om at sosiale effekter må følges opp gjennom bredere grasrot-initiativ og lavterskelarenaer.

\paragraph{Strategisk refleksjon}
Erfaringene fra Trondheim vil trolig prege Skiforbundets strategi:
\begin{itemize}
    \item Økt satsing på bærekraft og teknologi som konkurransefortrinn internasjonalt.
    \item Større vilje til å ta projekteierrollen i partnerskap med kommuner og næringsliv.
    \item En tydeligere publikums- og rekrutteringsstrategi mellom mesterskapene for å omsette VM-effekten til flere aktive medlemmer.
\end{itemize}

Dermed har Ski-VM ikke bare vært en målestokk på NSF sin prosjektkompetanse -- det har også gitt retning for forbundets posisjon i norsk og internasjonal skisport de neste årene.
\section{Oppgave 2}
\subsection{Valg av arbeidsgiver}
Twoday er \enquote{(...) en internasjonal leverandør av forretningskritiske IT-løsninger i Norden.}
\parencite{Twoday.no}. Ifølge nettsiden har twoday mer enn 2700 ansatte og over 8000 kunder. Innenfor IT
tilbyr selskapet en variasjon av områder å jobbe med som blant annet: Data og AI, Software Enginering
og Digital Experiences. Selskapet tilbyr variasjon i mulige arbeidsoppgaver, og nettsiden viser at de
verdsetter både generalister og spesialister. Tilpasningsevne og drivkraft er kjerneverdier hos twoday. Twoday er
en arbeidsgiver vi ser på som attraktiv, og vi vil videre diskutere hvordan deres fremstilling som seg selv som 
arbeidsgiver er motiverende på oss med grunnlag i pensumlitteraturen. 

Teorien om psykologiske jobbkrav vektlegger læring og variasjon presenteres som krav for å være tilfreds på jobb \parencite[][s.120]{Teknologiledelse}. 
Kjerneverdiene knyttet til tilpasningsevne og drivkraft er i lys av dette svært viktige for å skape motivasjon hos ansatte. Dette er også motiverende for 
oss da vi ønsker en arbeidsgiver hvor vi kan fortsette å utvikle oss og lære slik vi er vant til fra studiehverdagen.
Satsing på personlig utvikling og muligheten til å bestemme selv står sentralt hos twoday: \enquote{Du har muligheten
til å påvirke egen hverdag, og får en bratt læringskurve og mange utfordringer. Vår ambisjon er å
skape stolte øyeblikk(...).} \parencite["Graduateprogrammet"]{Twoday.no}. I tillegg til å vektlegge læring, ser
vi her at twoday vektlegger autonomi. Som beskrevet i Self Determination Theory\parencite[][s.121]{Teknologiledelse},
autonomi og å jobbe av egen fri vilje sentralt for å skape indre motivasjon hos mennesker. Læring,
variasjon og å bestemme over eget arbeid er sentralt for selvrealisering-toppen av Maslows
behovspyramide \parencite[][s.116]{Teknologiledelse}.

Twoday prioriterer mennesker, samhold, og bærekraft. Dette fanges opp verdien hjerte:
“Vi bryr oss om våre ansatte, våre kunder og samfunnet. Vi er en varm virksomhet som er fylt
med mennesker som har det gøy og ler sammen. Vi tar imot folk slik de er, vi samarbeider og vi
blomstrer sammen.” \parencite["Om oss"]{Twoday.no}. Vi motiveres av muligheten til å kunne knytte sosiale bånd
og utvikle oss sammen med likesinnede. Sosiale behov er sentrale i Marlows behovspyramide, å dekke
dette behovet på arbeidsplassen kan være med på å øke trivsel generelt. Vi motiveres også av å delta
i noe større og se fruktene av vårt arbeid. Dette stemmer overens med kravet: “sammenheng unngår
fremmedgjøring” fra teorien om psykologiske jobbkrav \parencite[][s.120]{Teknologiledelse}. I tillegg til å ha fokus på samhold
mellom de ansatte, er samfunnsnytte, og det å gjøre hverdagen enklere for folk viktig hos
twoday\parencite["Graduateprogrammet"]{Twoday.no}. Dette er motiverende for oss, da vi ønsker å delta i noe
større, og utgjøre en forskjell med den jobben vi gjør.

Med vekt på motiverende faktorer som sosialt samhold, selvutvikling, læring, mestring og utfordringer er derfor
twoday en motiverende arbeidsgiver. Flere sentrale motiverende faktorer fra litteraturen, som vi også subjektivt verdsetter
er tydelig i fokus hos dem, noe som gjør dem til en motiverende og interessant arbeidsgiver, \parencite["Graduateprogrammet"]{Twoday.no}. Dette er motiverende for oss, da vi ønsker å delta i noe
større, og utgjøre en forskjell med den jobben vi gjør.

Med vekt på motiverende faktorer som sosialt samhold, selvutvikling, læring, mestring og utfordringer er derfor
twoday en motiverende arbeidsgiver. Flere sentrale motiverende faktorer fra litteraturen, som vi også subjektivt verdsetter
er tydelig i fokus hos dem, noe som gjør dem til en motiverende og interessant arbeidsgiver. 




\section{Oppgave 3}
Den norske arbeidslivsmodellen blir kalt en triangelmodell ettersom den er bygd opp av tre aspekter, nemlig økonomisk styring, organisert arbeidsliv og offentlig velferd \cite{Teknologiledelse} (s.214). 
Videre kan den norske arbeidslivmodellen bli delt opp i tre nivåer - velferdsstatsmodellen, arbeidslivsmodellen og samarbeidsmodellen. 
Velferdsstatsmodellen handler om ordninger i arbeidslivet, derav arbeidsledighetstrygd, rett til ferie og sykeordninger. 
Dette er et aspekter ved den norske arbeidslivsmodellen som er viktig å ivareta. Dette er fordi ordninger som dette skaper forutsigbarhet i arbeidsmarkedet samt bidrar til økonomisk trygget for arbeidstakere.
 
Innenfor samarbeidsmodellen er gruppe- og teamarbeid svært sentral av den norske modellen. 
Teamarbeid handler om å utvikle team som får så mye ansvar som mulig gitt de rammebetingelsene man operer innenfor, og står team sentralt innenfor den norske modellen og er blitt sentralt i norsk arbeidsliv. 
Det å la teamet selv finne gode løsninger er en måte å ivareta tradisjonene i norsk arbeidsliv \cite{Teknologiledelse} (s.277), og er derfor et viktig aspekter ved den norske arbeidslivsmodellen som er viktig å ivareta. 
Dette er fordi teamarbeid fremmer selvstendighet og ansvarsfølelse, og når et team team gis frihet til å organisere sitt eget arbeid og finne løsninger innenfor gitte rammer, styrkes både innovasjonsevnen og effektiviteten i organisasjonen.
 
Ifølge en NRK-artikkel publisert januar 2025, er dagens generasjon mer krevende å lede. Den såkalte generasjon Z (Gen Z) mener derimot at det er sjefene som må endre metoder \cite{NRK}. Gen Z vil finne en balanse mellom jobb og fritid, spesielt muligheten til å jobbe mer digitalt og uavhengig av et fast kontor.  Mangelen på fleksibilitet, både når det gjelder arbeidstid og arbeidssted er en utfordring ved dagens arbeidsliv i Norge som bør endres for å bedre tilpasse seg Gen Z sine forventninger. Vi har et godt utgangspunkt for endring siden den norske arbeidslivsmodellen bygger på en triangellmodell mellom arbeidsgivere, arbeidstakere og myndigheter, samt høy grad av tillit og medbestemmelse på arbeidsplassen. Partnerskap mellom tillitsvalgte og ledelse i virksomheten er et viktig moment innenfor den norske modellen \parencite[229]{Teknologiledelse}. 
En mulig løsning er dermed å bruke tillitsvalgte for å inkludere unge ansatte (Gen Z) i beslutninger om fleksible arbeidsordninger. Dette sikrer at endringer skjer i tråd med både arbeidstakers og arbeidsgivers behov.
 
 
Et viktig særtrekk i norsk arbeidsliv er den sentrale rollen som fagforeninger har \parencite[214]{Teknologiledelse}. En annen mulig løsning er dermed å oppdatere tariffavtaler som åpner for større individuell fleksibilitet, for eksempel gjennom avtaler om hjemmekontor uten at det går på bekostning av arbeidsmiljø eller rettigheter. 

\section*{Konklusion}
Målet med dette prosjektet er å få en helhetlig innsikt i sentrale teamer innen teknologiledelse, prosjektarbeid og organisasjonsutvkling. 
Vi har anvendt teorier og modeller fra pensum for å analysere et konkret prosjekt, nemlig Ski-VM i Trondheim, og også utforsket personlige motivasjoner for karrierevalg samt studere hvordan det norske arbeidslivet kan utviklet i møte med Gen Z. 
Disse tre hovedoppgavene kobler teori og praksis sammen, og har gitt oss en bedre forståelse av både prosjektledelse og strategisk tenkning rundt fremtidens arbeidsliv. Gjennom dette prosjektet har vi styrket våre evner til kritisk refleksjon, både prosjektledelse og strategisk tenkning rundt fremtidens arbeidsliv. 

\printbibliography[heading=bibintoc, title={Referanser}]

\end{document}